\documentclass[../../../include/open-logic-section]{subfiles}

\begin{document}

\olfileid[fr]{sth}{z}{union}
\olsection{Union}

La \olref[sth][z][sep]{prop:intersectionsexist} nous a donné les
intersections. Mais si nous voulons que les unions arbitraires existent,
il nous faut poser un autre axiome :

\begin{axiom}[Union] Pour tout ensemble $A$, l'ensemble $\bigcup A =
\Setabs{x}{(\exists b \in A) x \in b}$ existe.
\[
	\forall A \exists U \forall x(x \in U \liff (\exists b \in A)x \in b)
\]
\end{axiom}

Cet axiome aussi se justifie par la conception cumulative et
itérative. Soit $A$ un ensemble : $A$ est donc formé à une certaine
étape $S$ (par \stageshier). Chaque !!{element} de $A$ a été formé
\emph{avant} $S$ (par \stagesacc) ; par un raisonnement semblable,
chaque !!{element} de chaque !!{element} de $A$ a donc été formé
avant $S$. Tous \emph{ces} ensembles sont ainsi disponibles avant
$S$, pour être réunis en un ensemble à l'étape $S$. Cet ensemble
n'est autre que $\bigcup A$.

\end{document}
